\newpage
\thispagestyle{empty}
\phantomsection
\addcontentsline{toc}{chapter}{Liste des abréviations}
\markboth{Liste des abréviations}{Liste des abréviations}

\begin{center}
\vspace*{0.5cm}
\begin{tikzpicture}
  \node[
    draw        = primary,
    line width  = 0.8pt,
    fill        = white,
    inner xsep  = 2.0cm,
    inner ysep  = 0.7cm,
    text width  = 10cm,
    align       = center
  ]{
    {\fontsize{16}{20}\selectfont\bfseries\scshape\rmfamily\color{primary}
    Liste des abréviations}
  };
\end{tikzpicture}
\vspace{0.8cm}
\end{center}

\renewcommand{\arraystretch}{1.45}
\arrayrulecolor{gray!40}
\rowcolors{2}{gray!5}{white}
\begin{longtable}{>{\bfseries\color{primary}}p{3.2cm} p{11.5cm}}

\rowcolor{blue!12}
\textbf{\color{black}Abréviation} & \textbf{\color{black}Signification} \\
\hline
\endfirsthead
\rowcolor{blue!12}
\textbf{\color{black}Abréviation} & \textbf{\color{black}Signification} \\
\hline
\endhead
\hline
\endfoot
\endlastfoot

% --- A ---
API   & \textit{Application Programming Interface} — Interface de programmation applicative \\
\hline

% --- B ---
B2B   & \textit{Business-to-Business} — Échanges commerciaux entre entreprises \\
\hline

% --- C ---
CI/CD & \textit{Continuous Integration / Continuous Deployment} — Intégration et déploiement continus \\
CORS  & \textit{Cross-Origin Resource Sharing} — Partage de ressources entre origines croisées \\
CSRF  & \textit{Cross-Site Request Forgery} — Falsification de requête inter-sites \\
CSV   & \textit{Comma-Separated Values} — Format de fichier à valeurs séparées par des virgules \\
\hline

% --- D ---
DSI   & Développement des Systèmes d'Information \\
\hline

% --- G ---
GLPI  & \textit{Gestionnaire Libre de Parc Informatique} \\
\hline

% --- H ---
HTTP  & \textit{HyperText Transfer Protocol} — Protocole de transfert hypertexte \\
HTTPS & \textit{HyperText Transfer Protocol Secure} — Version sécurisée du HTTP \\
\hline

% --- I ---
IA    & Intelligence Artificielle \\
IDOR  & \textit{Insecure Direct Object Reference} — Référence directe non sécurisée à un objet \\
ISET  & Institut Supérieur des Études Technologiques \\
ITSM  & \textit{Information Technology Service Management} — Gestion des services informatiques \\
\hline

% --- J ---
JSON  & \textit{JavaScript Object Notation} — Format d'échange de données léger \\
JWT   & \textit{JSON Web Token} — Jeton d'authentification au format JSON \\
\hline

% --- K ---
KB    & \textit{Knowledge Base} — Base de connaissances \\
KPI   & \textit{Key Performance Indicator} — Indicateur clé de performance \\
\hline

% --- L ---
LLM   & \textit{Large Language Model} — Grand modèle de langage \\
LPU   & \textit{Language Processing Unit} — Unité de traitement du langage (processeur Groq) \\
\hline

% --- M ---
MD5   & \textit{Message Digest 5} — Algorithme de hachage cryptographique \\
MIME  & \textit{Multipurpose Internet Mail Extensions} — Extensions multiusages pour le courrier Internet \\
MVC   & \textit{Model-View-Controller} — Patron de conception Modèle-Vue-Contrôleur \\
\hline

% --- O ---
OAuth & \textit{Open Authorization} — Protocole d'autorisation ouverte \\
ORM   & \textit{Object-Relational Mapping} — Mapping objet-relationnel \\
OWASP & \textit{Open Worldwide Application Security Project} — Projet mondial de sécurité applicative \\
\hline

% --- P ---
PDF   & \textit{Portable Document Format} — Format de document portable \\
PFE   & Projet de Fin d'Études \\
PII   & \textit{Personally Identifiable Information} — Informations personnelles identifiables \\
PV    & Procès-Verbal \\
\hline

% --- R ---
RBAC  & \textit{Role-Based Access Control} — Contrôle d'accès basé sur les rôles \\
REST  & \textit{Representational State Transfer} — Style d'architecture pour les services web \\
RPO   & \textit{Recovery Point Objective} — Objectif de point de reprise \\
RTO   & \textit{Recovery Time Objective} — Objectif de délai de reprise \\
\hline

% --- S ---
SaaS  & \textit{Software as a Service} — Logiciel en tant que service \\
SGBD  & Système de Gestion de Bases de Données \\
SLA   & \textit{Service Level Agreement} — Accord de niveau de service \\
SOLID & Acronyme des cinq principes de conception : \textit{Single responsibility, Open/closed, Liskov substitution, Interface segregation, Dependency inversion} \\
SQL   & \textit{Structured Query Language} — Langage de requête structuré \\
SSR   & \textit{Server-Side Rendering} — Rendu côté serveur \\
\hline

% --- T ---
TCP   & \textit{Transmission Control Protocol} — Protocole de contrôle de transmission \\
TLS   & \textit{Transport Layer Security} — Protocole de sécurisation des échanges réseau \\
TTC   & Toutes Taxes Comprises \\
TTL   & \textit{Time To Live} — Durée de vie (d'un cache ou d'un jeton) \\
TVA   & Taxe sur la Valeur Ajoutée \\
\hline

% --- U ---
UML   & \textit{Unified Modeling Language} — Langage de modélisation unifié \\
URL   & \textit{Uniform Resource Locator} — Localisateur uniforme de ressource \\
\hline

% --- W ---
WCAG  & \textit{Web Content Accessibility Guidelines} — Règles d'accessibilité du contenu web \\
\hline

% --- X ---
XP    & \textit{Extreme Programming} — Méthode agile de développement logiciel \\
XSS   & \textit{Cross-Site Scripting} — Injection de scripts malveillants côté client \\
\hline

% --- Y ---
YAML  & \textit{YAML Ain't Markup Language} — Format de sérialisation de données \\
\hline

% --- Z ---
ZAP   & \textit{Zed Attack Proxy} — Outil de test de sécurité web de l'OWASP \\
\hline

\end{longtable}
